% Generated draft: requires mathematical and source review.
\documentclass[11pt]{article}
\usepackage[margin=25mm]{geometry}
\usepackage{fontspec}
\setmainfont{Latin Modern Roman}
\usepackage{amsmath,amssymb,mathtools,unicode-math}
\setmathfont{Latin Modern Math}
\usepackage{enumitem,needspace,xurl,hyperref}
\setlength{\parindent}{0pt}
\setlength{\parskip}{5pt}
\setlength{\emergencystretch}{3em}
\newcommand{\N}{\mathbb{N}}
\newcommand{\Z}{\mathbb{Z}}
\newcommand{\Q}{\mathbb{Q}}
\newcommand{\R}{\mathbb{R}}
\newcommand{\C}{\mathbb{C}}
\newcommand{\F}{\mathbb{F}}
\newcommand{\A}{\mathbb{A}}
\newcommand{\PP}{\mathbb{P}}
\newcommand{\E}{\mathbb{E}}
\newcommand{\eps}{\varepsilon}
\newcommand{\dd}{\,\mathrm{d}}
\begin{document}
\section*{UnsolvedMath ID: 700345}
\textbf{Source problem code:} catalogue entry
\subsection{Definitions and mathematical statement}
A group \(G\) is \emph{periodic} if every element has finite order:
\[
\forall g\in G\;\exists n\in\Z_{\geq 1}\quad g^n=1.
\]
It is \emph{residually finite} if, for every \(g\in G\setminus\{1\}\), there exist a finite group \(F\) and a homomorphism \(\varphi:G\to F\) such that \(\varphi(g)\neq 1\).

For an explicit conditional formulation, interpret the \emph{weak minimum condition for subgroups} as the absence of an infinite descending chain
\[
H_0>H_1>H_2>\cdots,\qquad H_i\leq G,
\]
with \([H_i:H_{i+1}]=\infty\) for every integer \(i\geq0\). The supplied record does not define this condition; identification of this interpretation with the source's intended meaning requires review.

Under this interpretation, the question proposed by V. P. Shunkov is: must every periodic, residually finite group satisfying the weak minimum condition for subgroups be finite? Equivalently, does every infinite periodic, residually finite group admit an infinite descending chain of subgroups in which each successive index is infinite?
\subsection{Short English statement}
Must a periodic, residually finite group satisfying the weak minimum condition for subgroups be finite? The explicit chain interpretation above is conditional because the supplied record does not define that condition.
\subsection{Sources}
\begin{itemize}
\item Supplied local UnsolvedMath registry record, permanent ID \(700345\), exact source problem code \texttt{catalogue entry}, title \emph{Kourovka Notebook Problem 5.67}. The registry labels the problem \texttt{open}; this is supplied source metadata, not an independently verified status assessment.
\item V. P. Shunkov, problem 5.67, \emph{The Kourovka Notebook (Unsolved Problems in Group Theory)}, 5th issue (1976), as attributed in the supplied record.
\item The supplied background identifies the 21st issue edition (2026), arXiv:1401.0300v48: \url{https://arxiv.org/abs/1401.0300v48}. The supplied collection description names E. I. Khukhro and V. D. Mazurov as editors. This reference has not been independently consulted or verified here.
\end{itemize}
\end{document}
